\chapter{Conclusion de l'action}

\section{Bilan général}
La phase préparatoire nous a permis de construire une proposition fondée sur une démarche de diagnostic plutôt que sur un thème imposé. Après la prise de contact avec l'AFPA, nous avons complété les attentes du commanditaire par un questionnaire destiné aux stagiaires. L'analyse de treize réponses fait apparaître des lacunes importantes concernant l'alcool, ainsi que des comportements déclarés préoccupants liés au téléphone et aux distracteurs.

Cette préparation a également nécessité d'anticiper les conditions d'animation, de rechercher des ressources auprès de la Préfecture et de concevoir des situations adaptées aux déplacements du métier de plombier. À ce stade, ces acquis concernent la \emph{préparation} de l'action et non les résultats d'une sensibilisation déjà conduite.

\texteguide{Après l'intervention : compléter la conclusion avec le bilan réel, l'atteinte des objectifs, les limites du dispositif, les enseignements personnels en tant que futur ECSR et les perspectives de remédiation.}
